% !TEX root = ../main.tex
\chapter{Special Functions}
\label{ch:special_functions}

The mathematical special functions --- error functions, gamma
functions, Bessel functions, Legendre and Hermite polynomials,
and a handful of others --- appear in ML in specific contexts:
the error function inside GELU, the gamma function inside
log-likelihoods of Gamma and Beta distributions, the digamma in
variational inference, the Bessel function in Kaiser windows.
\Lucid\ exposes them at \code{lucid.special}, mirroring SciPy's
\code{scipy.special} where the names align.

This chapter catalogues the special-function family, the
numerical-precision considerations, and the implementation paths.

\section{Overview}
\label{sec:special_functions:overview}

\code{lucid.special} contains roughly 38 functions across
several families:

\begin{itemize}
  \item \textbf{Error functions}: \code{erf}, \code{erfc},
    \code{erfinv}, \code{erfcinv}.
  \item \textbf{Gamma family}: \code{lgamma}, \code{digamma},
    \code{trigamma}, \code{polygamma}.
  \item \textbf{Bessel family}: \code{i0}, \code{i1},
    \code{i0e}, \code{i1e}, \code{j0}, \code{j1}, \code{k0},
    \code{k1}, \code{y0}, \code{y1}.
  \item \textbf{Beta family}: \code{betainc}, \code{betaincinv}.
  \item \textbf{Orthogonal polynomials}: \code{legendre},
    \code{hermite}, \code{laguerre}, \code{chebyshev}.
  \item \textbf{Miscellaneous}: \code{expit} (sigmoid), \code{logit},
    \code{xlog1py}, \code{xlogy}, \code{exp2}, \code{expm1},
    \code{log1p}.
\end{itemize}

Each is a differentiable elementwise function from a real or
complex input to a real or complex output of the same shape.

\subsection{Why a separate subpackage}
\label{ssec:special_functions:why_separate}

The functions could in principle live alongside the elementwise
unary ops in \code{lucid.*}, but they are grouped in a subpackage
for parity with NumPy / SciPy / \refframework, all of which
expose them under \code{special} or similar. The grouping is
organisational; the implementations dispatch through the same
\code{UnaryKernel} machinery as everything else.

\section{Error Functions}
\label{sec:special_functions:error_funcs}

The error function $\operatorname{erf}(x) = \frac{2}{\sqrt{\pi}}
\int_0^x e^{-t^2} dt$ is the workhorse of the special-function
family in ML, primarily through GELU (\chref{ch:unary_functions})
and Gaussian quantile estimation.

\subsection{erf and erfc}
\label{ssec:special_functions:erf_erfc}

\code{erf(x)} returns $\operatorname{erf}(x) \in (-1, 1)$.
Backward: $2 e^{-x^2} / \sqrt{\pi}$.

\code{erfc(x) = 1 - erf(x)} is the complementary error function.
Equivalent mathematically but numerically more precise for
large $x$ where \code{1 - erf(x)} catastrophically cancels.
Backward: $-2 e^{-x^2} / \sqrt{\pi}$.

\subsection{erfinv and erfcinv}
\label{ssec:special_functions:erfinv}

The inverse error function. \code{erfinv(y)} returns $x$ such
that $\operatorname{erf}(x) = y$, for $y \in (-1, 1)$.

Implementation: \Accelerate's \code{vverfinvf} on CPU, an \MLX\
primitive on GPU. Both produce correctly rounded results in the
interior; near $|y| = 1$ the function grows unboundedly and
ULP-level precision is hard to maintain.

Backward: $\sqrt{\pi} \cdot e^{\operatorname{erfinv}(y)^2} / 2$.

\subsection{Use in Gaussian sampling}
\label{ssec:special_functions:gauss_sampling}

The inverse-CDF method for Gaussian sampling uses \code{erfinv}:
to sample from $\mathcal{N}(0, 1)$, draw $u \sim
\mathcal{U}(0, 1)$ and return $\sqrt{2} \cdot
\operatorname{erfinv}(2u - 1)$. \Lucid's \code{lucid.randn} uses
this path on CPU but a faster Ziggurat algorithm on GPU; the
\code{erfinv} method is the fallback.

\section{Gamma Family}
\label{sec:special_functions:gamma}

The gamma function $\Gamma(x)$ extends the factorial to real
arguments: $\Gamma(n + 1) = n!$ for non-negative integer $n$.
Its log $\log \Gamma$ is what \Lucid\ exposes, since the
absolute values of $\Gamma$ exceed FP32 range for arguments above
$\sim 35$.

\subsection{lgamma}
\label{ssec:special_functions:lgamma}

\code{lgamma(x)} returns $\log |\Gamma(x)|$. Implementation:
Stirling's approximation with correction terms; matches
\code{libm.dylib}'s \code{lgammaf} on CPU and \MLX's primitive
on GPU.

Backward: \code{digamma(x)}, the derivative of $\log \Gamma$.

\subsection{digamma and polygamma}
\label{ssec:special_functions:digamma}

\code{digamma(x)} returns $\psi(x) = \Gamma'(x) / \Gamma(x)$.
The digamma is the gradient of \code{lgamma}; \code{trigamma} is
the gradient of \code{digamma}; \code{polygamma(n, x)} is the
$n$-th derivative.

These appear in variational inference (the gradient of
Beta and Dirichlet log-likelihoods) and in Bayesian neural
networks.

\subsection{Why \texttt{lgamma}, not \texttt{gamma}}
\label{ssec:special_functions:why_lgamma}

\Lucid\ does not expose a bare \code{gamma}. The reason is range:
$\Gamma(35) \approx 10^{38}$, just barely below FP32's
$\sim 10^{38.5}$ overflow; $\Gamma(36)$ overflows. In practice,
ML almost always wants $\log \Gamma$ (the log-likelihood), so
exposing the log form by default avoids the overflow.

For users who genuinely want $\Gamma$, \code{lgamma(x).exp()}
gives it, with the same overflow risk.

\section{Bessel Functions}
\label{sec:special_functions:bessel}

Bessel functions arise in cylindrical PDEs and in some
window-function constructions (Kaiser uses the modified Bessel
$I_0$). \Lucid\ provides:

\begin{itemize}
  \item \code{i0(x), i1(x)}: modified Bessel of the first kind,
    orders 0 and 1.
  \item \code{i0e(x), i1e(x)}: exponentially-scaled variants,
    $I_n(x) e^{-|x|}$, useful for large $|x|$ where the
    unscaled version overflows.
  \item \code{j0(x), j1(x)}: ordinary Bessel of the first kind.
  \item \code{k0(x), k1(x)}: modified Bessel of the second kind.
  \item \code{y0(x), y1(x)}: ordinary Bessel of the second kind.
\end{itemize}

These are implemented via series expansion for small $|x|$ and
asymptotic forms for large $|x|$, with the crossover chosen for
precision continuity.

Backward: the Bessel functions satisfy recurrence relations that
express derivatives in terms of neighbouring orders
($I_0'(x) = I_1(x)$, etc.), so backward is straightforward.

\section{Orthogonal Polynomials}
\label{sec:special_functions:orthogonal}

Legendre, Hermite, Laguerre, and Chebyshev polynomials.
\code{legendre(n, x)} returns the $n$-th Legendre polynomial
evaluated at $x$.

These appear in spectral methods, in basis-function expansions
of probability distributions, and in physics-informed neural
networks. \Lucid's implementations use the standard recurrence
relations.

\section{Miscellaneous}
\label{sec:special_functions:misc}

\subsection{expit and logit}
\label{ssec:special_functions:expit_logit}

\code{expit(x) = 1 / (1 + exp(-x))} is sigmoid; we expose it
under both names for SciPy parity. \code{logit(p) = log(p / (1 -
p))} is its inverse.

The composite \code{lucid.expit} matches \code{lucid.sigmoid}; the
two are aliases.

\subsection{xlog1py and xlogy}
\label{ssec:special_functions:xlog_variants}

\code{xlogy(x, y) = x * log(y)}, with the convention that $0
\cdot \log 0 = 0$ (instead of NaN). Useful in entropy-style
computations where some entries may be zero.

\code{xlog1py(x, y) = x * log1p(y)} is the analogous variant for
$\log(1 + y)$.

These are typical entropy-loss helpers; the standard
\code{cross\_entropy} implementation uses them under the hood.

\subsection{log1p and expm1}
\label{ssec:special_functions:log1p_expm1}

Re-exported from the unary chapter; appear here for SciPy parity.
The implementations are the same.

\section{Numerical Precision}
\label{sec:special_functions:precision}

Special functions are where numerical precision matters most.
We summarise the precision of \Lucid's implementations.

\subsection{Precision goals}
\label{ssec:special_functions:precision_goals}

The target is ULP-level precision in FP32 across the input
range. Achieved for: \code{erf}, \code{erfc}, \code{lgamma},
\code{digamma}, \code{j0}, \code{j1}, \code{i0}, \code{i1},
\code{expit}, \code{logit}, \code{log1p}, \code{expm1}.

Achieved within a few ULPs (loose): \code{erfinv} near $|y| =
1$, \code{trigamma} for negative arguments, the higher-order
Bessel functions for very large arguments.

\subsection{The vForce vs scalar discrepancy}
\label{ssec:special_functions:vforce_scalar}

\Accelerate's vForce primitives are usually 1 ULP from the
scalar libm equivalents, but for \code{vverfinvf} the
discrepancy is larger (typically 4--6 ULPs near $|y| = 1$). The
test suite uses a relaxed tolerance for \code{erfinv}-using
ops.

This is documented in
\chref{ch:unary_functions}~Section~\ref{ssec:unary_functions:erfinv_precision}.

\section{Backward Derivatives}
\label{sec:special_functions:backward}

Most special functions have well-known derivatives:

\begin{table}[t]
\centering
\small
\begin{tabular}{ll}
\toprule
Function & Derivative \\
\midrule
$\operatorname{erf}(x)$    & $2 e^{-x^2} / \sqrt{\pi}$ \\
$\operatorname{erfc}(x)$   & $-2 e^{-x^2} / \sqrt{\pi}$ \\
$\operatorname{erfinv}(y)$ & $\sqrt{\pi} e^{\operatorname{erfinv}(y)^2} / 2$ \\
$\log \Gamma(x)$           & $\psi(x)$ (digamma) \\
$\psi(x)$                  & $\psi'(x)$ (trigamma) \\
$I_0(x)$                   & $I_1(x)$ \\
$I_1(x)$                   & $I_0(x) - I_1(x) / x$ \\
$P_n(x)$ (Legendre)        & via recurrence \\
\bottomrule
\end{tabular}
\caption[Special-function derivatives.]{Analytical derivatives
for the main special functions. \Lucid's autograd registers each
of these in the corresponding op's \code{BackwardSpec}.}
\label{tab:special_functions:derivatives}
\end{table}

\subsection{Higher-order derivatives}
\label{ssec:special_functions:higher_order}

The polygamma functions form a natural higher-order family: the
$n$-th derivative of $\log \Gamma$. \Lucid\ exposes \code{digamma}
($n=1$) and \code{trigamma} ($n=2$) explicitly; \code{polygamma(n, x)}
covers higher orders.

For an autograd chain that needs the second derivative of
\code{lgamma}, the engine composes \code{digamma} and
\code{trigamma} automatically.

\section{Use Cases in ML}
\label{sec:special_functions:use_cases}

\subsection{GELU activation}
\label{ssec:special_functions:gelu_uc}

\code{gelu(x) = x * (1 + erf(x / sqrt(2))) / 2}. The error
function is the integral of the standard normal density;
\code{gelu} is the expectation of $x \cdot \mathbb{1}[Z \leq
x]$ for $Z \sim \mathcal{N}(0, 1)$, a smoothed alternative to
ReLU that has become the default in modern Transformers.

\subsection{Gaussian quantile}
\label{ssec:special_functions:gauss_quantile}

The inverse CDF of $\mathcal{N}(0, 1)$ is $\sqrt{2}
\operatorname{erfinv}(2u - 1)$. Used in inverse-CDF sampling
(\chref{ch:distributions}).

\subsection{Beta and Dirichlet likelihoods}
\label{ssec:special_functions:beta_lik}

The Beta and Dirichlet densities involve $\Gamma$:
\[
p(x; \alpha, \beta) = \frac{\Gamma(\alpha + \beta)}{\Gamma(\alpha)
\Gamma(\beta)} x^{\alpha - 1} (1 - x)^{\beta - 1}.
\]
The log-likelihood uses \code{lgamma}; its gradient with respect
to $\alpha$ involves \code{digamma}.

\subsection{Kaiser window}
\label{ssec:special_functions:kaiser_uc}

The Kaiser window's shape comes from $I_0$, the modified Bessel
of order zero. \Lucid's \code{kaiser} window calls \code{i0}
directly.

\section{Implementation Notes}
\label{sec:special_functions:impl}

The implementations live in \code{lucid/special/} (Python) and
in \code{lucid/\_C/ops/ufunc/} (C++ where the op is a primitive).

\subsection{Composite vs primitive}
\label{ssec:special_functions:composite_vs_prim}

\code{erf}, \code{erfc}, \code{lgamma}, \code{digamma}, \code{i0},
\code{i1} are primitives, because they appear in hot paths
(GELU, distribution log-likelihoods).

\code{erfinv}, \code{trigamma}, \code{polygamma}, the higher-order
Bessel and Legendre functions are composites or thin wrappers
around primitives.

\subsection{Dispatch}
\label{ssec:special_functions:dispatch}

Special-function primitives dispatch through the standard
backend mechanism (\chref{ch:backend_dispatch}). The CPU path
goes through \Accelerate's \code{vForce}; the GPU path uses
\MLX's primitives. The user does not see the difference.

\section{The Gamma Function in Depth}
\label{sec:special_functions:gamma_depth}

The gamma function $\Gamma$ generalises the factorial:
\[
\Gamma(z) \;=\; \int_0^\infty t^{z-1} e^{-t}\, dt,
\quad \Re(z) > 0.
\]
Its key properties:

\begin{itemize}
  \item Recursion: $\Gamma(z + 1) = z\,\Gamma(z)$, hence
    $\Gamma(n+1) = n!$ for non-negative integer $n$.
  \item Reflection: $\Gamma(z)\Gamma(1-z) = \pi /
    \sin(\pi z)$.
  \item Stirling asymptotic: $\Gamma(z) \sim \sqrt{2\pi/z}\,
    (z/e)^z$ for large $|z|$.
  \item Special values: $\Gamma(1/2) = \sqrt{\pi}$,
    $\Gamma(3/2) = \sqrt{\pi}/2$.
\end{itemize}

\subsection{The log-gamma identity}
\label{ssec:special_functions:lgamma_identity}

A practical identity for ML log-likelihoods, the
\emph{log-Beta} function:
\[
\log \mathrm{B}(\alpha, \beta) \;=\; \log\Gamma(\alpha) +
\log\Gamma(\beta) - \log\Gamma(\alpha + \beta).
\]
This appears in the log-likelihood of the Beta distribution and
in the Dirichlet's normalisation constant. \Lucid's \code{lgamma}
implementation is what these depend on.

\subsection{The digamma and harmonic numbers}
\label{ssec:special_functions:digamma_harmonic}

The digamma function $\psi(z) = \Gamma'(z) / \Gamma(z)$ relates
to the harmonic numbers:
\[
\psi(n+1) \;=\; -\gamma + \sum_{k=1}^n \frac{1}{k},
\]
where $\gamma \approx 0.5772$ is the Euler--Mascheroni constant.
At half-integer arguments,
\[
\psi(1/2) \;=\; -\gamma - 2\ln 2.
\]
These identities appear in entropy formulas for Gamma and Beta
distributions.

\subsection{Reflection identity in practice}
\label{ssec:special_functions:reflection}

The reflection identity lets the implementation compute $\Gamma$
for negative arguments via positive ones:
\[
\Gamma(-z) \;=\; \frac{-\pi}{z\,\Gamma(z)\,\sin(\pi z)}
\quad (z > 0,\ z \notin \mathbb{Z}).
\]
\Lucid's \code{lgamma} uses this internally; the user sees the
unified surface.

\section{The Error Function in Depth}
\label{sec:special_functions:erf_depth}

The error function
\[
\operatorname{erf}(x) \;=\; \frac{2}{\sqrt{\pi}} \int_0^x e^{-t^2}\, dt
\]
is the cumulative distribution of a half-Gaussian, scaled.
Properties:

\begin{itemize}
  \item Odd: $\operatorname{erf}(-x) = -\operatorname{erf}(x)$.
  \item Bounds: $-1 < \operatorname{erf}(x) < 1$.
  \item Derivative: $\operatorname{erf}'(x) = 2 e^{-x^2} / \sqrt{\pi}$.
  \item Asymptotic: $\operatorname{erf}(x) \to 1$ as $x \to
    \infty$, with $1 - \operatorname{erf}(x) \sim e^{-x^2} /
    (x\sqrt{\pi})$.
\end{itemize}

\subsection{The standard normal CDF}
\label{ssec:special_functions:phi_via_erf}

The CDF of the standard normal $\mathcal{N}(0, 1)$ is
\[
\Phi(x) \;=\; \frac{1}{2}\bigl(1 + \operatorname{erf}(x/\sqrt{2})\bigr).
\]
This identity is what makes GELU's smooth-Gaussian form
expressible as
\[
\mathrm{GELU}(x) \;=\; x \cdot \Phi(x)
\;=\; \frac{x}{2}\bigl(1 + \operatorname{erf}(x/\sqrt{2})\bigr).
\]

\subsection{The complementary error function for tails}
\label{ssec:special_functions:erfc_tails}

For $x \gtrsim 4$ in FP32, $1 - \operatorname{erf}(x)$ loses
precision through cancellation. The complementary form
$\operatorname{erfc}(x) = 1 - \operatorname{erf}(x)$ is
implemented directly:
\[
\operatorname{erfc}(x) \;=\; \frac{2}{\sqrt{\pi}} \int_x^\infty e^{-t^2}\, dt.
\]
For Gaussian tail probabilities ($P(Z > z)$ for $Z \sim
\mathcal{N}(0, 1)$), the recommended formula is
\[
P(Z > z) \;=\; \frac{1}{2} \operatorname{erfc}(z / \sqrt{2}),
\]
which preserves precision even at $z \gtrsim 5$.

\section{Bessel Functions in Depth}
\label{sec:special_functions:bessel_depth}

The modified Bessel function of the first kind, order $\nu$, is
the function $I_\nu$ satisfying
\[
x^2 y'' + x\, y' - (x^2 + \nu^2) y \;=\; 0.
\]
For $\nu = 0$:
\[
I_0(x) \;=\; \sum_{k=0}^\infty \frac{(x/2)^{2k}}{(k!)^2}
        \;=\; \frac{1}{\pi} \int_0^\pi e^{x \cos\theta}\, d\theta.
\]

\subsection{Asymptotic forms}
\label{ssec:special_functions:bessel_asymptotic}

For $x \to \infty$,
\[
I_\nu(x) \;\sim\; \frac{e^x}{\sqrt{2\pi x}}
\;\left(1 - \frac{4\nu^2 - 1}{8x} + O(x^{-2})\right).
\]
The $e^x$ factor causes overflow in FP32 for $x \gtrsim 90$. The
exponentially-scaled variant
\[
I_\nu^{e}(x) \;=\; I_\nu(x) \cdot e^{-|x|}
\]
removes the exponential and is finite throughout. \Lucid\
exposes \code{i0e}, \code{i1e} for these scaled forms.

\subsection{Recurrence relations}
\label{ssec:special_functions:bessel_recurrence}

The orders satisfy
\[
I_{\nu - 1}(x) - I_{\nu + 1}(x) \;=\; \frac{2\nu}{x} I_\nu(x),
\]
and derivatives:
\[
I_\nu'(x) \;=\; \tfrac{1}{2}\bigl(I_{\nu - 1}(x) + I_{\nu + 1}(x)\bigr).
\]
The recurrence allows computing higher orders from lower ones;
\Lucid\ exposes only $\nu = 0, 1$ directly (the most common in
applications) and the user composes higher orders if needed.

\section{Orthogonal Polynomials in Depth}
\label{sec:special_functions:orthogonal_depth}

Each orthogonal polynomial family satisfies a three-term
recurrence that makes evaluation efficient.

\subsection{Legendre polynomials}
\label{ssec:special_functions:legendre_detail}

Defined on $[-1, 1]$ with weight $w(x) = 1$:
\[
P_0(x) = 1,\quad P_1(x) = x,
\]
\[
(n+1) P_{n+1}(x) \;=\; (2n+1) x P_n(x) - n P_{n-1}(x).
\]
Orthogonality:
\[
\int_{-1}^{1} P_m(x) P_n(x)\, dx \;=\; \frac{2}{2n+1} \delta_{mn}.
\]

\subsection{Hermite polynomials}
\label{ssec:special_functions:hermite_detail}

The probabilist's form on $\mathbb{R}$ with weight $w(x) =
e^{-x^2/2}$:
\[
He_0(x) = 1,\quad He_1(x) = x,
\]
\[
He_{n+1}(x) \;=\; x\, He_n(x) - n\, He_{n-1}(x).
\]
These appear in Hermite-Gauss quadrature and in series
expansions of Gaussian processes.

\subsection{Chebyshev polynomials}
\label{ssec:special_functions:chebyshev_detail}

First kind, on $[-1, 1]$ with weight $w(x) = 1/\sqrt{1 - x^2}$:
\[
T_0(x) = 1,\quad T_1(x) = x,
\]
\[
T_{n+1}(x) \;=\; 2 x T_n(x) - T_{n-1}(x).
\]
Chebyshev polynomials are optimal in the minimax sense for
polynomial approximation; they underlie the Chebyshev iteration
in iterative linear solvers.

\section{The Beta Function and Beta Distribution}
\label{sec:special_functions:beta_function}

The Beta function
\[
\mathrm{B}(\alpha, \beta) \;=\; \int_0^1 t^{\alpha - 1} (1 - t)^{\beta - 1}\, dt
\;=\; \frac{\Gamma(\alpha) \Gamma(\beta)}{\Gamma(\alpha + \beta)}.
\]
appears in the normalisation of the Beta distribution
\[
p(x; \alpha, \beta) \;=\; \frac{x^{\alpha - 1} (1 - x)^{\beta - 1}}{\mathrm{B}(\alpha, \beta)},
\quad x \in (0, 1).
\]
The Beta function's log form uses \code{lgamma}:
\[
\log \mathrm{B}(\alpha, \beta) \;=\; \log\Gamma(\alpha) + \log\Gamma(\beta) - \log\Gamma(\alpha + \beta).
\]
The Beta log-likelihood gradient with respect to $\alpha$ is
\[
\frac{\partial}{\partial \alpha} \log p(x; \alpha, \beta)
\;=\; \log x \;-\; \psi(\alpha) \;+\; \psi(\alpha + \beta),
\]
involving the digamma function.

\subsection{The incomplete Beta}
\label{ssec:special_functions:betainc}

\code{betainc(a, b, x)} returns the regularised incomplete Beta
function
\[
I_x(a, b) \;=\; \frac{1}{\mathrm{B}(a, b)} \int_0^x t^{a-1} (1-t)^{b-1}\, dt,
\]
which is the CDF of the Beta distribution. \code{betaincinv} is
the inverse, used in confidence-interval and quantile
computations.

\section{Numerical Precision Summary}
\label{sec:special_functions:precision_summary}

\tabref{tab:special_functions:precision_summary} summarises
\Lucid's per-function precision targets.

\begin{table}[t]
\centering
\small
\begin{tabular}{lll}
\toprule
Function & FP32 precision & Notes \\
\midrule
\code{erf}, \code{erfc}    & $\leq 2$ ULP & throughout $\mathbb{R}$ \\
\code{erfinv}              & $\leq 5$ ULP & loose near $|y| = 1$ \\
\code{lgamma}              & $\leq 3$ ULP & all $x > 0$; reflection for $x < 0$ \\
\code{digamma}             & $\leq 4$ ULP & loose near poles \\
\code{i0}, \code{i1}       & $\leq 3$ ULP & full range with $i0e/i1e$ for large $|x|$ \\
\code{j0}, \code{j1}       & $\leq 5$ ULP & looser due to oscillation \\
\code{legendre(n, x)}      & $\leq n + 2$ ULP & recurrence amplifies error \\
\code{betainc}             & $\leq 6$ ULP & convergence-dependent \\
\bottomrule
\end{tabular}
\caption[Special-function precision.]{Approximate precision
targets for \Lucid's special-function implementations. The
oscillating Bessel functions and the high-order Legendre
polynomials are looser; the dominant entropy-likelihood family
(\code{lgamma}, \code{digamma}, \code{erf}) is tight to a few
ULPs.}
\label{tab:special_functions:precision_summary}
\end{table}

\section{Composite vs Primitive Cost}
\label{sec:special_functions:composite_cost}

For the ops implemented as composites (orthogonal polynomials
via recurrence, higher-order Bessel via downward recurrence),
the per-call cost is dominated by the recurrence depth. Measured
per-call costs on M3~Max for $10^6$-element FP32 tensors:

\begin{table}[t]
\centering
\small
\begin{tabular}{lrr}
\toprule
Function & Path & Time (ms) \\
\midrule
\code{erf}                & primitive (MLX/vForce) & 0.6 \\
\code{lgamma}             & primitive              & 1.1 \\
\code{i0}                 & primitive              & 1.4 \\
\code{legendre(8, x)}     & composite (8 recurrences) & 6.7 \\
\code{legendre(64, x)}    & composite (64 recurrences) & 52 \\
\code{betainc(2, 3, x)}   & composite              & 15 \\
\bottomrule
\end{tabular}
\caption[Special-function timings.]{Per-call timings for
representative special functions on M3~Max with $10^6$-element
FP32 inputs. The primitive ops are bandwidth-bound; the
composite ops scale with the recurrence depth.}
\label{tab:special_functions:timings}
\end{table}

\section{Summary and Forward References}
\label{sec:special_functions:summary}

\code{lucid.special} provides the standard family of special
functions used in ML: error functions, gamma family, Bessel
family, orthogonal polynomials. The implementations target
ULP-level FP32 precision where possible and document the few
cases (\code{erfinv} near $|y|=1$, higher-order Bessels at large
arguments) where the precision is more loose.

The next chapter, \chref{ch:distributions}, treats the
probability-distribution subpackage that uses several of these
functions internally.

\begin{lucidnote}
For the user: \code{lucid.special} is the namespace; most ops are
called as \code{lucid.special.erf(x)}, \code{lucid.special.lgamma(x)},
and similar. The autograd machinery treats each as a primitive,
with the analytical derivatives registered.
\end{lucidnote}
